\chapter{Finite canonical support}\label{ch-support}\label{sec-finite}

Lemma~\ref{contact-structure} identifies a maximising graphon almost
everywhere with the binary threshold kernel $\mathbf1_{\{F(x)\cdot F(y)>0\}}$
built from its own canonical frame. The canonical law $\nu=F_\#dx$ is a
priori an arbitrary Borel probability measure on the ellipsoid
$\{v:v^TLv=\mathcal M\}$, so the kernel could still be supported on a
continuum of frame values. This chapter proves that it is not: every maximiser
has a canonical law with at most $24$ atoms, all of positive mass. That is
Theorem~\ref{structure-main}, whose proof is assembled here from four steps.

The first step, Theorem~\ref{negative}, is a second-variation inequality.
Multiplying the underlying measure by $1+tg(F(x))$ with $\int g\,d\nu=0$ is a
legitimate change of graphon for small $|t|$, and maximality forces the
second-order term to be nonpositive. The resulting kernel
$K(v,w)=(v\cdot w)_+- (v\cdot w)(v^TLw)$ is negative semidefinite \emph{on the
actual support}; no extension to the whole sphere is claimed or needed.

The second step, Lemma~\ref{propagate}, converts null vectors of that kernel
into vector identities. On a set of mutually positively correlated points, a
relation $\sum_i a_iv_iv_i^T=0$ forces both $\sum_ia_iv_i=0$ and
$\sum_ia_iA(v_i,w)=0$ for every $w$. Since the space of symmetric $3\times3$
matrices has dimension six and the sphere splits into four sign regions on
which all dot products are positive, this bounds the dimension of the row
space of $A$ by $24$ and provides actual anchor points $z_j\in S$.

The third step, Lemma~\ref{descent} and Theorem~\ref{finite}, turns the
dimension bound into a cardinality bound. Inside a single sign region of the
anchor equators the rows $A(v,\cdot)$ are linear functions of $v$, so
Lemma~\ref{descent} applies and gives at most $2^{d-1}\le4$ points per region;
the number of regions of $r$ lines is at most $1+r(r+1)/2$. This yields the
intermediate bound $|S|\le1204$.

The fourth step, Section~\ref{sec-separation}, improves $1204$ to $24$ by
excluding minimal second-moment dependencies on positive cliques.
Lemma~\ref{binary} is an elementary circuit argument that uses a separating
family of binary observations; the actual rows $B_{ki}$ of the maximiser
supply such a family. The conclusion, Theorem~\ref{six}, is that on every
positive clique the map $a\mapsto\sum_{i\in J}a_iv_iv_i^T$ is injective, so a
clique has at most six types; four sign regions then give
$|S|\le4\cdot6=24$.

\section{A negative semidefinite kernel on the actual support}
% Ported verbatim from paper/top3_joint_moment_bound.tex lines 444-487.
% Only cross-reference wording was adapted; no mathematics was altered.
Continue with the maximizing graphon and canonical law of the previous
section, and write $h_v(w)=\mathbf1_{\{v\cdot w>0\}}$. Define continuous symmetric kernels on $S\times S$ by
\[
 A(v,w)=(v\cdot w)_+,\qquad
 D(v,w)=(v\cdot w)(v^TLw),\qquad K=A-D.
\]
\begin{theorem}[Second variation on the actual support]\label{negative}
For every finite list $v_1,\ldots,v_s\in S$ and every real list $a_1,\ldots,a_s$,
\begin{equation}
 \sum_{i,j}a_i a_j K(v_i,v_j)\le0. \label{nsd}
\end{equation}
The assertion concerns only the actual canonical support.
\end{theorem}
\begin{proof}
For any bounded measurable real $g$ on $S$ with $\int g\,d\nu=0$, use the legitimate density
$1+tg(F(x))$ for both signs of sufficiently small $t$.
Set
\[
 Q=\int g(v)vv^T\,d\nu(v),\qquad
 R=\iint h_v(w)vw^Tg(v)g(w)\,d\nu(v)d\nu(w).
\]
The unchanged three coordinate functions have the complete Gram $G_t=I+tQ$
and numerator matrix $N_t=L+t(QL+LQ)+t^2R$ in this changed measure.
For small $t$, $G_t$ is positive definite; multiplying the frame by $G_t^{-1/2}$
is a genuine common orthonormalization. Since
$N_t=G_tLG_t+t^2(R-QLQ)$ and $\tr(LQ)=\mathcal M\int g=0$, its exact trace is
\begin{equation}
 \tr(G_t^{-1}N_t)=\mathcal M+t^2\tr\bigl(G_t^{-1}(R-QLQ)\bigr)\le \mathcal M. \label{density}
\end{equation}
Letting $t\to0$ gives $\iint K(v,w)g(v)g(w)\,d\nu(v)d\nu(w)\le0$.
There is no discarded remainder or varying-scheme constant in this calculation.
By Lemma~\ref{ellipse-structure} and~\eqref{all-rows}, $\int K(v,w)d\nu(w)=\mathcal M-\mathcal M=0$ for every $v\in S$.
Subtracting the mean therefore proves the same inequality for every bounded $g$.
To obtain \eqref{nsd}, combine repeated points and take disjoint shrinking
neighborhoods of the remaining points. Each has positive $\nu$-mass by full support.
Use $g=\sum_i a_i1_{U_i}/\nu(U_i)$, separately for each neighborhood size.
This auxiliary $g$ is bounded but need not have mean zero. The legal probability
density direction is $g-\int g\,d\nu$; centering leaves the $K$ quadratic form unchanged.
Its bound need not be uniform across neighborhood sizes.
Continuity of $K$ on compact $S\times S$ gives \eqref{nsd} in the limit.
This limiting step needs no uniform permissible size of $t$.
No eigenvalue differentiation is used, including when $\lambda_2=\lambda_3$.
\end{proof}


\section{Propagation of moment relations}\label{sec-propagation}
% Ported verbatim from paper/top3_joint_moment_bound.tex lines 489-519.
% Only cross-reference wording was adapted; no mathematics was altered.
\begin{lemma}\label{propagate}
If $v_i\in S$ have $v_i\cdot v_j>0$ for all $i,j$, then
\[
 \sum_i a_i v_iv_i^T=0
 \quad\Longrightarrow\quad
 \sum_i a_i v_i=0\quad\hbox{and}\quad
 \sum_i a_i A(v_i,w)=0\quad\hbox{for every }w\in S.
\]
Consequently the vector space spanned by the functions $A(v,\cdot)$, $v\in S$,
has dimension $r\le24$ and admits a basis $A(z_j,\cdot)$ with actual $z_j\in S$.
\end{lemma}
\begin{proof}
Write $Q_a=\sum a_i v_iv_i^T$ and $b_a=\sum a_i v_i$.
The quadratic form in \eqref{nsd} is $|b_a|^2-\tr(LQ_a^2)$ on this positive-dot set.
For $Q_a=0$ it is nonnegative, hence zero and $b_a=0$.
A null vector of a negative semidefinite quadratic form has zero cross terms:
apply \eqref{nsd} also to this list with $w$ appended and coefficient $s\in\R$.
The resulting polynomial $2s\sum_i a_iK(v_i,w)+s^2K(w,w)$ is nonpositive for all $s$,
so its linear coefficient vanishes. The $D$ sum vanishes because $Q_a=0$;
therefore the $A$ sum vanishes for every $w$.
Partition $S$ by the signs of coordinates 2 and 3, assigning zero to the positive sign.
Within each of the four disjoint sets, every dot product is strictly positive,
since coordinate 1 is positive and the other two coordinate products are nonnegative.
In each nonempty set choose a basis, of size at most 6, of its matrices $vv^T$
in the six-dimensional space of symmetric $3$ by $3$ matrices.
The implication just proved expresses every corresponding $A$ row in the span
of those basis-point rows. Combine the four lists and discard dependent rows.
This proves $r\le24$. These are spanning witnesses, not a replacement vertex measure.
This bounds the row space, not yet the cardinality of the support.
\end{proof}


\section{Sign regions and finite support}\label{sec-signregions}
% Ported verbatim from paper/top3_joint_moment_bound.tex lines 521-577.
% Only cross-reference wording was adapted; no mathematics was altered.
\begin{lemma}[Sign-region dimension descent]\label{descent}
Let $T\subset S$ have linear span of dimension $d\ge1$.
Suppose for every $w\in S$ there is a vector $b_w\in\R^3$ such that
$A(v,w)=v\cdot b_w$ for all $v\in T$. Then $|T|\le2^{d-1}$.
\end{lemma}
\begin{proof}
If $d=1$, the positive first coordinate and $v^TLv=\mathcal M$ make $T$ a singleton.
For general $d$, if no $w\in S$ has strictly positive dot product with one member
of $T$ and strictly negative dot product with another, any two members have the
same threshold row outside the union of their two zero equators.
Equation \eqref{all-rows} then makes the members equal.
Otherwise choose such a $w$, with witnesses $v_+,v_-$, and split
$T_+=\{v\in T:v\cdot w\ge0\}$ and $T_-=\{v\in T:v\cdot w<0\}$.
On $T_+$, $v\cdot(b_w-w)=0$; on $T_-$, $v\cdot b_w=0$.
Both are proper hyperplanes of $\operatorname{span}T$: evaluate $b_w-w$ at $v_-$
and $b_w$ at $v_+$. Thus both children have span dimension at most $d-1$.
They inherit the linearization for every $w'\in S$.
Induction gives at most $2^{d-2}$ points in each child, proving the bound.
Zeros were included in $T_+$, so none are lost in the dimension descent.
\end{proof}

\begin{theorem}[Finite support of every maximizer]\label{finite}
For every global maximizer of $\Sigma_3$, its actual canonical support satisfies
\begin{equation}
 |S|\le 4\bigl(1+r(r+1)/2\bigr)\le1204=:N_*. \label{bound}
\end{equation}
The original graphon agrees almost everywhere with a binary step graphon
on these actual fibers, with
positive masses, diagonal blocks 1, and no equal rows. For this particular
maximizer $\min_{v,w\in S}|v\cdot w|>0$.
\end{theorem}
\begin{proof}
Take the actual anchors $z_1,\ldots,z_r$ from Lemma~\ref{propagate}.
Symmetry gives, for each $w\in S$, coefficients $c_j(w)$ with
$A(v,w)=\sum_j c_j(w)A(v,z_j)$ for every $v\in S$.
The union of the finitely many anchor equators has $\nu$-measure zero by \eqref{all-rows}.
On any fixed strict sign region of these anchors,
$A(v,z_j)$ is either $v\cdot z_j$ or zero. Thus $A(v,w)$ is linear in $v$
there, for every $w$, with coefficients determined by that region.
Because $v_1>0$, use coordinates $(v_2/v_1,v_3/v_1)$.
Each anchor equator becomes an affine line, or is empty if its last two coordinates vanish.
At most $r$ lines have at most $1+r(r+1)/2$ complementary regions: the $j$-th
line has at most $j-1$ distinct previous intersection points and adds at most $j$ regions.
Coincident or parallel lines only lower this count. A strict sign pattern is an
intersection of open half-planes, hence connected whenever nonempty.
Apply Lemma~\ref{descent} to $T=S$ intersected with each such sign region;
its span dimension is at most 3, so it contains at most four points.
We have therefore placed full $\nu$-measure in a finite set of size bounded by \eqref{bound}.
Its closure is the same finite set, so full support of $\nu$ implies that all of $S$
is contained in it, including any previously excluded boundary candidates.
Every support point is now an atom of positive mass. The individual zero-equator
statements in \eqref{all-rows} imply $v\cdot w\ne0$ for every pair of support atoms.
Taking a minimum over this finite set proves the asserted margin.
The actual fibers $F^{-1}(v)$ provide the partition, and \eqref{all-rows} identifies its rows.
\end{proof}
The margin asserted here may depend on the maximizer.


\section{Separating binary observations}\label{sec-separation}
% Ported verbatim from paper/top3_joint_moment_bound.tex lines 579-620.
% Only cross-reference wording was adapted; no mathematics was altered.
By the preceding section, write $S=\{v_1,\ldots,v_m\}$ with actual masses
$p_i>0$. Put $V=(v_i^T)_{i=1}^m$, $P_0=\diag(p_i)$ and
$B_{ij}=\mathbf1_{\{v_i\cdot v_j>0\}}$. Then
\begin{equation}\label{finite-rows}
 BP_0V=VL,\qquad \sum_i p_iv_iv_i^T=I_3,\qquad
 v_i^TLv_i=\mathcal M.
\end{equation}
The full rows of $B$ are distinct: equal rows would imply equal $v_i$
by invertibility of $L$. A positive clique $J$ means a subset of this actual support with $B_{ij}=1$
for every $i,j\in J$. Extend vectors indexed by $J$ by zero outside $J$.
\begin{lemma}[Vector rigidity on a positive clique]
For every real $a$ supported on a positive clique, if
$Q_a=\sum_i a_i v_iv_i^T=0$, then
\begin{equation}
 \sum_i a_i=0,\quad b_a:=\sum_i a_iv_i=0,\quad
 \sum_i B_{ki}a_iv_i=0\quad\hbox{for every row }k. \label{rigid}
\end{equation}
\end{lemma}
\begin{proof}
Taking the trace against $L$ gives $\mathcal M\sum a_i=0$.
For both signs of sufficiently small nonzero $t$, all $p_i+ta_i$ are positive.
These are legitimate new fiber measures for the same binary template, realizable
on an atomless space. They are not independent choices of moment data.
The original coordinate functions have Gram matrix $I+tQ_a=I$ and compression
\[
 N(t)=L+t(Q_aL+LQ_a)+t^2b_ab_a^T=L+t^2b_ab_a^T.
\]
The last term uses the clique assumption, and both cross terms use \eqref{finite-rows}.
Their trace cannot exceed the global maximum $\mathcal M$; hence $b_a=0$ and $N(t)=L$.
Equality of this trial trace with the maximum forces its three-dimensional
space to be invariant. Explicitly, rotate one orthonormal coordinate into any
unit perpendicular function; the first derivative of the trace gives zero
off-space matrix element. This argument does not differentiate ordered eigenvalues.
The compression $L$ then gives the eigenfunction equations at the perturbed measure.
All fibers still have positive mass, so for every $k$,
$\sum_i(p_i+ta_i)B_{ki}v_i=Lv_k$.
Subtract \eqref{finite-rows} and divide by $t$ to obtain \eqref{rigid}.
\end{proof}
Moment-preserving cloning and the Ky Fan equality argument follow the
mechanism in~\cite[Claims 2.2--2.5]{KMMP}. The permissible size of $t$
may depend on the maximizer and the chosen direction.


\subsection{A separating observation lemma}
% Ported verbatim from paper/top3_joint_moment_bound.tex lines 622-716.
% Only cross-reference wording was adapted; no mathematics was altered.
Minimal dependencies of rank-one powers have a broader theory; see
Li and Liu~\cite{LiLiu}. The following lemma uses the additional binary
separation hypothesis. Its elementary proof is included.
\begin{lemma}[Moment independence from separating binary observations]\label{binary}
Let $T=\{v_i:i\in J\}$ be a finite set of distinct points in $\R^3$ with
$v_{i1}>0$ and $v_i^TLv_i=\mathcal M>0$, where $L$ is positive definite.
Let $\mathcal B$ be a family of functions $b:J\to\{0,1\}$ containing the
constant-one function and separating every pair of distinct indices.
Suppose, for every real $a$ with $\sum_{i\in J}a_iv_iv_i^T=0$, that
\begin{equation}
 \sum_{i\in J}a_i b(i)v_i=0\quad\hbox{for every }b\in\mathcal B. \label{obs}
\end{equation}
Then the matrices $\{v_iv_i^T:i\in J\}$ are linearly independent.
\end{lemma}
\begin{proof}
Suppose instead that $I\subseteq J$ is a support-minimal dependent family
of these matrices, with relation $a_i\ne0$ for all $i\in I$.
Write $s=|I|$. Since $\dim\operatorname{Sym}_3=6$, minimality gives $s\le7$.
The ellipse equation gives $\sum_{I}a_i=0$, while the constant observation
in \eqref{obs} gives $\sum_I a_iv_i=0$.
With $w_i=(1,v_i)^T\in\R^4$, all three identities combine to give
\[
 \sum_{a_i>0}a_iw_iw_i^T=\sum_{a_i<0}(-a_i)w_iw_i^T.
\]
Both sides are positive semidefinite. The range of a positive weighted sum
of $w_iw_i^T$ is precisely the span of those $w_i$:
its kernel consists of vectors perpendicular to every such $w_i$.
Consequently the two spans coincide, and the affine hull $H$ of all $v_i$
in the circuit satisfies
\[
 1+\dim H=\dim\operatorname{span}\{w_i:i\in I\}
 \le\min(|\{a_i>0\}|,|\{a_i<0\}|)\le3.
\]
This is an elementary positive-Gram range argument, not a new general circuit bound.
Distinct points here cannot be proportional: equal ellipse radius forces
the proportionality factor to be $\pm1$, and the positive first coordinate
excludes minus one. Also an affine line meets this positive-definite ellipsoid
in at most two points, by its strictly quadratic restriction to the line.
If $\dim H\le1$, there are at most two points, whose rank-one matrices
are independent by nonproportionality, a contradiction.
It remains that $\dim H=2$.

If $0\in H$, this is a two-dimensional linear plane. Its rank-one matrices
span a space of dimension at most three, so minimality now gives $s\le4$.
For any observation $b$ nonconstant on $I$, \eqref{obs} and the total vector
relation give nonzero-coefficient linear relations on both nonempty sets
$I_1=\{i:b(i)=1\}$ and $I_0=\{i:b(i)=0\}$.
Nonproportionality forces each set to contain at least three points.
This would give $s\ge6$, a contradiction. If instead every observation in
$\mathcal B$ were constant on $I$, then minimality excludes $s=1$ (a
one-element circuit is impossible because $v_iv_i^T\ne0$), so we may choose
distinct $i,j\in I$; since $\mathcal B$ separates every pair of distinct
indices of $J$, some $b\in\mathcal B$ satisfies $b(i)\ne b(j)$,
contradicting constancy on $I$.

If $0\notin H$, write $H=\{v:c\cdot v=1\}$.
Any three distinct points of $H$ on the ellipsoid are linearly independent:
a linear relation has coefficient sum zero after applying $c$, hence would
make the three points affinely collinear, which the line-intersection fact forbids.
A nonconstant observation would therefore require at least four points
on each of $I_1,I_0$, yielding $s\ge8$, again a contradiction.
Thus every observation is constant on $I$. A one-element circuit is impossible
because $v_iv_i^T\ne0$; separation of any two indices in $I$ is contradicted.
\end{proof}

\begin{theorem}[Every actual positive clique has at most six types]\label{six}
For every actual maximizer, and every positive
clique $J$ of its full support, the map
$a\mapsto\sum_{i\in J}a_iv_iv_i^T$ is injective. In particular $|J|\le6$.
If $\{v_i:i\in J\}$ lies in a linear plane, then $|J|\le3$.
If it lies in an affine plane not containing zero, then $|J|\le5$.
\end{theorem}
\begin{proof}
Use the actual observations $b_k(i)=B_{ki}$, for all rows $k$.
For nonempty $J$, a row $k\in J$ supplies the constant-one observation.
They separate indices: $B$ is symmetric and has distinct full rows.
Lemma \ref{binary} now applies by \eqref{rigid}.
The linear-plane refinement follows from $\dim\operatorname{Sym}_2=3$.
For an affine plane $c\cdot v=1$, the nonzero symmetric matrix
$L-\mathcal M cc^T$ is orthogonal to every $v_iv_i^T$. It is nonzero because $L$
has rank three while $\mathcal M cc^T$ has rank at most one. This leaves dimension at most five.
Empty cliques require no argument.
\end{proof}

\begin{corollary}[Proof of Theorem~\ref{structure-main}]\label{twentyfour}
Every maximizer has at most 24 distinct canonical support points.
\end{corollary}
\begin{proof}
Partition its full support by the signs of coordinates two and three,
assigning a zero coordinate to the positive sign. There are at most four parts.
Within each part, the last two coordinate products are nonnegative and the
first-coordinate product is strictly positive. Each part is a positive clique.
Theorem \ref{six} gives $m\le4\cdot6=24$ on its original support. This proves Theorem~\ref{structure-main}.
\end{proof}

